\documentclass[12pt]{article}
\usepackage{amsmath,amssymb,amsfonts}
\usepackage[margin=1in]{geometry}

\begin{document}

\section*{Chapter 2, Problem 8}

\subsection*{Problem Statement}
\textit{Brachistochrone problem for a central gravitational force.} A particle of mass $m$ is to move under the gravitational attraction of another stationary spherical mass $M$ from a point at a finite distance $r = R_0$ from the center of $M$ to a point on the surface of $M$. We want to determine the path such that the particle moving along it will travel between the two fixed endpoints in a time less than the time of transit along any other path. Taking for the two endpoints $\theta = 0$, $r = R_0$ and $\theta = \theta_s$,

$r = R$ (as shown in Fig.~2.7), set up in polar coordinates the equation that determines the path and solve for $\theta$ as a function of $r$. The best one can do analytically is to express $\theta$ as an elliptic integral with a variable upper limit and lower limit chosen to satisfy one of the boundary conditions. The determination of this path tells us how best to design interplanetary mail chutes.

\subsection*{Solution}

In polar coordinates centered on $M$, the line element is $d\ell^2 = dr^2 + r^2 d\theta^2$. By conservation of energy (the particle starts from rest at $r = R_0$):
\[
  \frac{1}{2}m v^2 - \frac{GMm}{r} = -\frac{GMm}{R_0},
\]
so
\[
  v^2 = 2GM\left(\frac{1}{r} - \frac{1}{R_0}\right).
\]

The travel time is
\[
  T = \int \frac{d\ell}{v} = \int \frac{\sqrt{dr^2 + r^2 d\theta^2}}{v}.
\]
Using $\theta$ as the independent variable with $r = r(\theta)$ and $r' = dr/d\theta$:
\[
  T = \int_0^{\theta_s} \frac{\sqrt{r'^2 + r^2}}{\sqrt{2GM(1/r - 1/R_0)}}\,d\theta.
\]
The integrand $f(\theta, r, r') = \sqrt{r'^2 + r^2}/\sqrt{2GM(1/r - 1/R_0)}$ does not depend explicitly on $\theta$, so the first integral applies:
\[
  r'\frac{\partial f}{\partial r'} - f = -C,
\]
where $C > 0$ is a constant. Computing:
\[
  \frac{\partial f}{\partial r'} = \frac{r'}{\sqrt{r'^2+r^2}\sqrt{2GM(1/r-1/R_0)}},
\]
\[
  r'\frac{\partial f}{\partial r'} - f = \frac{r'^2 - r'^2 - r^2}{\sqrt{r'^2+r^2}\sqrt{2GM(1/r-1/R_0)}} = \frac{-r^2}{\sqrt{r'^2+r^2}\sqrt{2GM(1/r-1/R_0)}}.
\]
So:
\[
  \frac{r^2}{\sqrt{r'^2+r^2}\sqrt{2GM(1/r-1/R_0)}} = C.
\]
Squaring:
\[
  \frac{r^4}{(r'^2+r^2)\cdot 2GM(1/r - 1/R_0)} = C^2.
\]
Solving for $r'^2$:
\[
  r'^2 + r^2 = \frac{r^4}{2GM\,C^2(1/r - 1/R_0)},
\]
\[
  r'^2 = \frac{r^4}{2GM\,C^2(1/r - 1/R_0)} - r^2
  = r^2\left[\frac{r^2}{2GM\,C^2(1/r - 1/R_0)} - 1\right].
\]

Let $\alpha = 2GM\,C^2$. Then:
\[
  r'^2 = r^2\left[\frac{r^2}{\alpha(1/r - 1/R_0)} - 1\right]
  = r^2\left[\frac{r^3}{\alpha(1 - r/R_0)} - 1\right].
\]

Separating variables:
\[
  d\theta = \frac{r'\,d\theta}{r'} = \frac{dr}{r'} = \frac{dr}{r\sqrt{\frac{r^3}{\alpha(1-r/R_0)} - 1}}.
\]

Therefore:
\[
  \theta = \int_{R}^{r} \frac{dr}{r\sqrt{\dfrac{r^3}{\alpha(1-r/R_0)} - 1}},
\]
where the limits are chosen so that $\theta = \theta_s$ when $r = R$ (the surface of $M$) and $\theta = 0$ when $r = R_0$.

Writing this more explicitly with the substitution cleaned up:
\[
  \theta(r) = \int_{R_0}^{r} \frac{dr'}{r'\sqrt{\dfrac{r'^3 R_0}{\alpha(R_0 - r')} - 1}}.
\]

This integral is an elliptic integral (it involves a square root of a cubic/quartic polynomial in $r'$). The constant $\alpha$ (equivalently $C$) is determined by the boundary condition $\theta(R) = \theta_s$.

In general, this integral cannot be evaluated in closed form and must be expressed as an elliptic integral with a variable upper limit, as stated in the problem.

\end{document}
